% !TEX program = lualatex
% Generated by new_lecture_note.sh
\documentclass[11pt]{article}

\usepackage{fontspec}
\usepackage{microtype}
\usepackage{mathtools,amssymb,amsfonts,amsthm,bm}
\usepackage{enumitem}
\usepackage{booktabs,array,xcolor}
\usepackage{geometry,fancyhdr,setspace}
\usepackage[hidelinks]{hyperref}
\usepackage[nameinlink,noabbrev]{cleveref}

\geometry{margin=1in,headheight=15pt}
\setstretch{1.12}
\setlength{\parindent}{0pt}
\setlength{\parskip}{0.55em plus 0.12em minus 0.08em}
\setlength{\abovedisplayskip}{0.8em plus 0.2em minus 0.1em}
\setlength{\belowdisplayskip}{0.8em plus 0.2em minus 0.1em}
\allowdisplaybreaks[3]
\setlist{leftmargin=*,topsep=0.35em,itemsep=0.15em,parsep=0pt}

% ---------- Metadata ----------
\newcommand{\studentname}{Keshav Ramamurthy}
\newcommand{\coursename}{Math 104}
\newcommand{\courselongname}{Real Analysis}
\newcommand{\lecturenumber}{02}
\newcommand{\lecturetitle}{Lecture 02}
\newcommand{\lecturedate}{\today}

% ---------- Reusable notation ----------
\newcommand{\N}{\mathbb{N}}
\newcommand{\Z}{\mathbb{Z}}
\newcommand{\Q}{\mathbb{Q}}
\newcommand{\R}{\mathbb{R}}
\newcommand{\C}{\mathbb{C}}
\newcommand{\F}{\mathbb{F}}
\newcommand{\K}{\mathbb{K}}
\newcommand{\E}{\mathbb{E}}
\newcommand{\Prob}{\mathbb{P}}
\newcommand{\1}{\mathbf{1}}
\newcommand{\eps}{\varepsilon}
\newcommand{\dd}{\,\mathrm{d}}
\newcommand{\abs}[1]{\left\lvert#1\right\rvert}
\newcommand{\norm}[1]{\left\lVert#1\right\rVert}
\newcommand{\ip}[2]{\left\langle#1,#2\right\rangle}
\newcommand{\set}[1]{\left\{#1\right\}}
\newcommand{\given}{\,\middle\vert\,}
\newcommand{\indep}{\perp\!\!\!\perp}
\newcommand{\lto}{\longrightarrow}
\newcommand{\deriv}[2]{\frac{\mathrm{d}#1}{\mathrm{d}#2}}
\newcommand{\pderiv}[2]{\frac{\partial#1}{\partial#2}}
\DeclareMathOperator{\Aut}{Aut}
\DeclareMathOperator{\Hom}{Hom}
\DeclareMathOperator{\Ker}{ker}
\DeclareMathOperator{\im}{im}
\DeclareMathOperator{\rank}{rank}
\DeclareMathOperator{\nullity}{nullity}
\DeclareMathOperator{\Span}{span}
\DeclareMathOperator{\Var}{Var}
\DeclareMathOperator{\Cov}{Cov}

% ---------- Note environments ----------
\newtheoremstyle{notestyle}{0.7em}{0.7em}{\itshape}{}{}{\bfseries}{.5em}{}
\theoremstyle{notestyle}
\newtheorem{theorem}{Theorem}[section]
\newtheorem{lemma}[theorem]{Lemma}
\newtheorem{proposition}[theorem]{Proposition}
\newtheorem{corollary}[theorem]{Corollary}
\theoremstyle{definition}
\newtheorem{definition}[theorem]{Definition}
\newtheorem{example}[theorem]{Example}
\newtheorem{remark}[theorem]{Remark}
\newenvironment{claim}{\par\noindent\textbf{Claim.}\ }{\par}
\newenvironment{idea}{\par\noindent\textit{Idea.}\ }{\par}
\newenvironment{proofsketch}{\par\noindent\textbf{Proof sketch.}\ }{\par}

% ---------- Header ----------
\pagestyle{fancy}
\fancyhf{}
\fancyhead[L]{\coursename}
\fancyhead[C]{\lecturetitle}
\fancyhead[R]{\studentname}
\fancyfoot[C]{\thepage}

\title{\vspace{-1.5em}\textbf{\coursename\ -- \lecturetitle}}
\author{\studentname}
\date{\lecturedate}

\begin{document}
\maketitle

\section{Doing $\epsilon-N$ proofs}
\begin{example}[1]
\label{ex:label}
Show that $\lim_{n\rightarrow\infty} \frac{1}{\sqrt{n}} = 0.$
\end{example}
\begin{proof}
    Square it, starting with $$|\frac{1}{\sqrt{n}} - 0| < \epsilon$$
    $$\frac{1}{n} < e^2 \rightarrow n > \frac{1}{e^2}$$
\end{proof}
You define $N = \frac{1}{e^2}$, and when $n > N$ the limit holds.
\begin{example}[2]
\label{ex:label}
Show that $\lim_{n\rightarrow\infty} \frac{3n+1}{7n+4} = \frac{3}{7} $

\end{example}
\begin{proof}
    We do this by writing the format first $$\left|\frac{3n+1}{7n+4} - \frac{3}{7}\right| < \epsilon$$
\end{proof}
We take a common denominator, to get $$\left|\frac{(3n+1)7 - 3(7n+4)}{(7n+4)(7)}\right|
< \epsilon$$
Then, this simplifes to $\left|\frac{19}{(7n-4)(7)}\right|$ which we can pull out absolute value as we see
$(7n-4)$ is positive.
The inequality then simplies to $\frac{19}{49\epsilon} + \frac{4}{7} < n$ and we set
$N = \frac{19}{49\epsilon}+\frac{4}{7}$
\par \textbf{Proof Format: } 
\par To do $\epsilon - N$ proofs, you just define $N = \text { something in terms of
\epsilon}.$ Then, set $n > N$, and show that the absolute value expression can then be
rearranged to show the original expression.
% \textit{Write the one-sentence point of today's lecture here.}
\begin{example}[3]
\label{ex:label}
Show that $\lim_{n\rightarrow\infty} \frac{4n^2 + 3n}{n^2 - 6}= 4$
\par\textbf{Scratch Work: } We want the inequality $$\left|\frac{4n^2+3n}{n^2-6}-4\right|<\epsilon$$
We take the common denominator:$$\left| \frac{(4n^2+3n)-4(4n^2-6)}{n^2-6}\right| $$
which is $$\left|\frac{3n^2+24}{n^2-6}\right| < \epsilon \rightarrow \text{ notice that
both numerator and denominator go to infinity}$$
So we can remove the absolute value.
Something here we do is we notice $n > 1$ so we write $$\frac{3n+24}{n^2-6}
<\frac{27n}{n^2-6} < \frac{27n}{\frac{1}{2}n^2} < \frac{54}{n}$$ for sufficiently large n.
\par So you just do $N = \max\{\frac{54}{\epsilon}, 12\}.$
% \frac{27n}{n^2-6}$$
\end{example}
\section{Definitions and notation}
\begin{definition}[First definition]
Write the definition and one sentence explaining why it matters.
\end{definition}

\section{Claims, theorems, and proof sketches}
\begin{claim}
State the claim in one precise sentence.
\end{claim}

\begin{proofsketch}
Record the key idea, the first useful equation, and the step that still needs checking.
\end{proofsketch}

\section{Examples and calculations}
\begin{example}[Lecture example]
Write the setup, the calculation, and the conclusion.
\end{example}

\section{Questions and follow-up}
\begin{itemize}
  \item What is still unclear?
  \item Which exercise or reading should I revisit?
  \item TODO:
\end{itemize}

\end{document}
